%% ============================================================
\section{Complete Formalization of CRAFT}
\label{app:formal-method}
%% ============================================================

This section expands the method in Section~\ref{sec:method} into a complete
input-to-output specification.  Given a task with hidden goals, CRAFT masks
those goals and applies the same rollout--decompose--pool--filter--compose
operator at inference time and during label construction.  We first define
the program semantics and verifier, then the exact and resource-bounded
filters, and finally the inference, SFT, and RL procedures.  Shapley credit and group reward collapse are formalized alongside RL scoring
and policy updates; the remaining theoretical properties and proofs are
collected in Appendix~\ref{app:composition-proofs} within this formalization.

\subsection{Task Semantics, Loop Invariants, and Verification}
\label{app:formal-loop-verification}

Let \(\Sigma\) be the program-state space.  A loop is
\(\mathcal L=(\mathsf{Init},B,S)\), where \(\mathsf{Init}\) characterizes the
first loop-head states, \(B\) is the guard, and
\(S\subseteq\Sigma\times\Sigma\) is one normally completing loop step.  Thus
\(S(\sigma,\sigma')\) includes execution of the body and any guard or update
side effects required to reach the next loop head.

\begin{definition}[Inductive invariant]
\label{def:formal-inductive}
A predicate \(I\) is a valid loop invariant when
\begin{equation}
  \begin{split}
  &\forall\sigma\in\Sigma:\
    \mathsf{Init}(\sigma)\Rightarrow I(\sigma),\\
  &\forall\sigma,\sigma'\in\Sigma:\
    I(\sigma)\land B(\sigma)\land S(\sigma,\sigma')
    \Rightarrow I(\sigma').
  \end{split}
  \label{eq:formal-inductive}
\end{equation}
We denote these two obligations by \(\operatorname{Ind}_{\mathcal L}(I)\).
Here and throughout the paper, \emph{valid invariant} means inductive in this
sense.  This implies validity on every reachable loop-head state, although the
converse need not hold for the chosen assertion language and prover.
\end{definition}

A response is decomposed into a finite candidate family \(C\) of clauses.  The
clauses compose by conjunction,
\begin{equation}
  I(C)=\bigwedge_{c\in C}c,
  \qquad I(\varnothing)=\mathit{true}.
  \label{eq:formal-compose}
\end{equation}
The set \(C\) is \emph{jointly inductive} when
\(\operatorname{Ind}_{\mathcal L}(I(C))\) holds.  Notice that preservation is
checked under the full conjunction: a clause may rely on companion clauses.
The logical definition is insensitive to clause order.  The bounded verifier,
however, receives a deterministic sequence
\(D=\langle d_1,\ldots,d_m\rangle\), because generated proof dependencies can
follow annotation order.  Set notation below denotes the members of this
sequence; deletion preserves their relative order.

A complete task is \(\mathcal T=(\mathcal L,\mathcal G)\), where
\(\mathcal G\) contains the assertions and postconditions at their original
program locations.  For a body target, let
\(S^{B}_{\ell}\subseteq\Sigma\times\Sigma\) be the path-sensitive prefix
relation from the loop head through guard evaluation to \(\ell\).
For a target after the loop, let \(E_0\) be the guard-false exit relation
and \(E_b\) the relation from a loop head through the body prefix to a
\texttt{break} that exits this loop.  These relations include guard and
branch conditions and all side effects.  Let \(F_{e,\ell}\) be the suffix
from exit \(e\) to \(\ell\), the identity for an immediately following
target.  Define the complete exit-to-target relation by
\[
 X_\ell(\sigma,\sigma')\iff
 \bigvee_{e\in\{0\}\cup\mathcal B}
 \exists\eta:\ E_e(\sigma,\eta)\land F_{e,\ell}(\eta,\sigma'),
\]
where \(\mathcal B\) indexes the loop's break exits.  For a side-effect-free
guard, \(E_0(\sigma,\eta)\) is \(\neg B(\sigma)\land\eta=\sigma\).
Unlike \(E_0\), a break exit may start with \(B(\sigma)\) true.
Then define
\begin{equation}
  \operatorname{VC}_{\mathcal L}(I,\ell,g)=
  \begin{cases}
    \forall\sigma,\sigma':
    I(\sigma)\land S^{B}_{\ell}(\sigma,\sigma')
    \Rightarrow g(\sigma'),
      & \ell\text{ is inside the body},\\
    \forall\sigma,\sigma':
    I(\sigma)\land X_\ell(\sigma,\sigma')
    \Rightarrow g(\sigma'),
      & \ell\text{ is after the loop}.
  \end{cases}
  \label{eq:formal-target-vc}
\end{equation}
Postconditions are represented by the second case, with suffix relations
including assignments and return-value substitution.  A return or other
abrupt exit inside the body, when admitted, requires its own exit relation
and applicable target obligations in the same union; it does not contribute
to the normally completing transition \(S\).  For example, with
\texttt{while (1) \{ if (q) break; ... \}}, the guard-false branch is empty,
but the break branch still yields a nonvacuous post-loop obligation.
The ideal logical verdict is
\begin{equation}
  \operatorname{TaskValid}_{\models}(\mathcal T,C)
  =\operatorname{Ind}_{\mathcal L}(I(C))\land
   \bigwedge_{(\ell,g)\in\mathcal G}
   \operatorname{VC}_{\mathcal L}(I(C),\ell,g).
  \label{eq:formal-verify}
\end{equation}
Generation and training see only \(\mathcal L\); \(\mathcal G\) is restored
only for Equation~\eqref{eq:formal-verify}.

For a proof goal \(g\) generated by the configured loop verifier, define its
portfolio status by
\begin{equation}
  \operatorname{GoalStatus}_{\delta}(g)
  \in\{\mathsf{valid},\mathsf{timeout},\mathsf{unknown}\},
  \label{eq:formal-goal-status}
\end{equation}
where \(\mathsf{valid}\) means that at least one configured prover returned
\emph{Valid}; \(\mathsf{timeout}\) means that no prover proved the goal and at
least one exhausted \(\delta\); and \(\mathsf{unknown}\) contains every other
non-valid, unsupported, missing, or malformed result.

Let \(\mathcal S_{\mathrm{ver}}=
\{\mathsf{valid},\mathsf{timeout},\mathsf{unknown}\}\).  One invocation
of the resource-bounded loop verifier on the current ordered family
\(D=\langle d_1,\ldots,d_m\rangle\) returns the clause-to-status map
\begin{equation}
  \operatorname{LoopVerify}_{\mathcal L,\delta}(D)
  =\{d_k\mapsto v_{D,k}:1\le k\le m\}
  \in\mathcal S_{\mathrm{ver}}^{D},
  \label{eq:formal-loop-verify-map}
\end{equation}
where \(v_{D,k}=\mathsf{valid}\) means that the verifier proves all
establishment and preservation obligations generated for \(d_k\) in the
current family \(D\).  If the obligations are not all proved,
\(v_{D,k}=\mathsf{timeout}\) records an exhausted prover budget and
\(v_{D,k}=\mathsf{unknown}\) records every other outcome.  The interface is
sound when
\begin{equation}
  \left(\forall c\in D:\
    \operatorname{LoopVerify}_{\mathcal L,\delta}(D)[c]=\mathsf{valid}\right)
  \Longrightarrow \operatorname{Ind}_{\mathcal L}(I(D)).
  \label{eq:formal-loop-verify-soundness}
\end{equation}
The filter treats \(\mathsf{timeout}\) and \(\mathsf{unknown}\) identically,
but the interface keeps them distinct.  Appendix~\ref{app:loop-verifier-realization}
describes how the implementation generates and aggregates the clause goals.

After the invariant set is fixed, the separate binary operator
\(\operatorname{TaskVerify}_{\delta}(\mathcal T,C)\) runs the verifier on the
original task annotated with \(C\).  Let
\(\operatorname{TargetGoals}(\mathcal T,C)\) be its generated source-level
assertion and postcondition goals.  Then
\begin{equation}
\begin{split}
\operatorname{TaskVerify}_{\delta}(\mathcal T,C)
=\mathbf 1\bigl[&\operatorname{SyntaxOK}(\mathcal T,C)\ \land\\
&\forall c\in C:\
  \operatorname{LoopVerify}_{\mathcal L,\delta}(C)[c]=\mathsf{valid}\ \land\\
&\forall g\in\operatorname{TargetGoals}(\mathcal T,C):
  \operatorname{GoalStatus}_{\delta}(g)=\mathsf{valid}\bigr].
\end{split}
\label{eq:formal-task-verify}
\end{equation}

\subsection{Ideal and Resource-Bounded Houdini Filtering}
\label{app:formal-houdini}

We define an exact semantic fixed-point operator and its resource-bounded
verifier counterpart.

For \(c\in C\), define its establishment and context-dependent preservation
obligations by
\begin{equation}
  \begin{split}
  \operatorname{Est}_{\mathcal L}(c)
    &\iff \forall\sigma:\
      \mathsf{Init}(\sigma)\Rightarrow c(\sigma),\\
  \operatorname{Pres}_{\mathcal L,C}(c)
    &\iff \forall\sigma,\sigma':
      I(C)(\sigma)\land B(\sigma)\land S(\sigma,\sigma')
      \Rightarrow c(\sigma').
  \end{split}
  \label{eq:formal-clause-obligations}
\end{equation}
Ideal Houdini starts with \(C_0=C\) and repeats
\begin{equation}
  C_{j+1}=\mathsf F_{\mathcal L}(C_j)
  =\{c\in C_j:\operatorname{Est}_{\mathcal L}(c)
    \land\operatorname{Pres}_{\mathcal L,C_j}(c)\}
  \label{eq:formal-houdini}
\end{equation}
until a fixed point \(\mathsf H_{\mathcal L}(C)\) is reached.

The implemented deletion round uses the clause-level loop verifier and
preserves the relative order of retained clauses:
\begin{equation}
  \widehat{\mathsf F}_{\mathcal L,\delta}(D)
  =\left\langle c\in D\ \middle|\
    \operatorname{LoopVerify}_{\mathcal L,\delta}(D)[c]=\mathsf{valid}
  \right\rangle.
  \label{eq:formal-bounded-houdini-round}
\end{equation}
Thus timeout and unknown both cause conservative deletion.  Every status in
one round is computed under the same pre-deletion context \(D\).  A clause
removed at the end of that round may therefore have helped another clause
receive \(\mathsf{valid}\).  The survivors must be verified again without the
removed hypotheses; this is why a single deletion pass is insufficient.  For
an ordered input sequence \(D\), let \(D_0=D\) and
\(D_{j+1}=\widehat{\mathsf F}_{\mathcal L,\delta}(D_j)\).  Iteration stops at
the first proved fixed point or at the symbolic round limit \(h_{\max}\).  Let
\begin{equation}
  j^*=\min\{j<h_{\max}:
    \widehat{\mathsf F}_{\mathcal L,\delta}(D_j)=D_j\},
  \label{eq:formal-bounded-fixed-point}
\end{equation}
when this set is nonempty.  The implemented filter is
\begin{equation}
  \operatorname{Filter}_{\mathcal L}^{\delta,h_{\max}}(D)=
  \begin{cases}
    D_{j^*}, & j^*\text{ exists},\\
    \langle\rangle, & \text{the round limit is reached first}.
  \end{cases}
  \label{eq:formal-bounded-houdini}
\end{equation}
At a returned fixed point, every surviving clause has just been proved in the
context of that same sequence.  Reaching the limit returns the empty sequence.  Under a
sound loop verifier, the result is sound and may be smaller than the ideal
fixed point.  In
the remainder, \(\operatorname{Filter}_{\mathcal L}\) abbreviates
Equation~\eqref{eq:formal-bounded-houdini}, and
\(\operatorname{TaskVerify}\) abbreviates
\(\operatorname{TaskVerify}_{\delta}\).

\subsection{The Shared Composition Operator}

Let \(\operatorname{parts}(y)\) extract the ordered sequence of invariant
clauses in response \(y\), let \(\operatorname{First}_{m}(L)\) be the prefix
of sequence \(L\) containing at most \(m\) clauses, and let \(\Vert\) denote
sequence concatenation.  Define the clauses admitted from one response by
\begin{equation}
  \operatorname{Admit}(y)=
  \operatorname{Unique}\!\left(
    \operatorname{First}_{m_{\max}}(\operatorname{parts}(y))
  \right).
  \label{eq:formal-admit}
\end{equation}
\(\operatorname{Unique}\) removes duplicates using the finite-rule
canonical normalization \(\operatorname{key}(c)\).  Equal keys identify the
clauses merged by those rules.  Provenance is retained by
recording which responses supplied each key.  It keeps the first occurrence;
concatenations are traversed in response order and then clause order, which
fixes the sequence supplied to \(\operatorname{LoopVerify}\).
For a clause sequence \(D\), let
\(\operatorname{keys}(D)=\langle\operatorname{key}(c):c\in D\rangle\).

For any response sequence \(y_{1:n}\), define
\begin{align}
  A_i&=\operatorname{Admit}(y_i),
  &P(y_{1:n})&=\operatorname{Unique}\!\left(A_1\Vert\cdots\Vert A_n\right),
  \label{eq:formal-pool}\\
  C(y_{1:n})&=\operatorname{Filter}_{\mathcal L}(P(y_{1:n})),
  &I(y_{1:n})&=\bigwedge_{c\in C(y_{1:n})}c.
  \label{eq:formal-craft-map}
\end{align}
This is the shared rollout--decompose--pool--filter--compose operator.  The
filter checks preservation under the complete current conjunction, so clauses
from different responses may support one another; standalone preservation is
not required.

\subsection{Complete Inference Algorithm}

\begin{algorithm}[H]
\caption{CRAFT inference}
\label{alg:rcf}
\begin{algorithmic}
\alginput{task \(\mathcal T=(\mathcal L,\mathcal G)\), policy \(\pi\), number of responses \(n\)}
\algoutput{composed invariant \(I\), final verdict \(v\)}
\algline \(x\leftarrow\operatorname{Mask}(\mathcal T)=\mathcal L\)\cmnt{hide goals}\\
\algline sample \(y_{1:n}\sim\pi(\cdot\mid x)\)\cmnt{rollout}\\
\algline \(A_i\leftarrow\operatorname{Admit}(y_i)\) for every \(i\)\cmnt{decompose and cap}\\
\algline \(P\leftarrow\operatorname{Unique}(A_1\Vert\cdots\Vert A_n)\)\cmnt{pool}\\
\algline \(C\leftarrow\operatorname{Filter}_{\mathcal L}^{\delta,h_{\max}}(P)\)\cmnt{filter}\\
\algline \(I\leftarrow\bigwedge_{c\in C}c\)\cmnt{compose}\\
\algline \(v\leftarrow\operatorname{TaskVerify}_{\delta}(\mathcal T,C)\)\cmnt{restore goals}\\
\algline \kw{return} \((I,v)\)\\
\end{algorithmic}
\end{algorithm}

For evaluation, compose@\(n\) is the percentage of tasks for which
\(\operatorname{TaskVerify}(\mathcal T,C(y_{1:n}))\) succeeds on the first
\(n\) saved responses. Pass@\(n\) is the percentage of tasks for which
at least one of these same \(n\) responses independently verifies the targets
(Appendix~\ref{app:experimental-protocol}).

\subsection{Complete SFT Label Construction}

SFT distills the same multi-response map into one response.  It accumulates
raw candidates so that a clause discarded in one round can later acquire a
supporting companion.  The policy \(\pi_0\) is the small open-weight model
being trained and is the only learned generator in label construction.
Executions of \(\mathcal L\) produce the trace signal, and
\(\operatorname{Filter}_{\mathcal L}\) supplies the verifier decision; no
external model supplies candidates, labels, or acceptance judgments.

\begin{algorithm}[H]
\caption{Multi-response SFT synthesis}
\label{alg:formal-sft}
\begin{algorithmic}
\alginput{masked loop \(\mathcal L\), policy \(\pi_0\), negatives \(\mathcal N\), group size \(g_{\mathrm{sft}}\), round limit \(t_{\mathrm{sft}}\)}
\algoutput{one serialized label, or \(\mathsf{reject}\)}
\algline \kw{if} \(\mathcal N=\varnothing\) \kw{then return} \(\mathsf{reject}\)\\
\algline \(P\leftarrow\langle\rangle\)\\
\algline \kw{for} \(t=1,\ldots,t_{\mathrm{sft}}\) \kw{do}\\
\algline \algind sample a response group \(y_{1:g_{\mathrm{sft}}}\sim\pi_0(\cdot\mid\mathcal L)\)\\
\algline \algind \(P\leftarrow\operatorname{Unique}(P\Vert
  \operatorname{Admit}(y_1)\Vert\cdots\Vert
  \operatorname{Admit}(y_{g_{\mathrm{sft}}}))\)\\
\algline \algind \(C\leftarrow\operatorname{Filter}_{\mathcal L}(P)\)\\
\algline \algind \kw{if} \(\operatorname{cov}_{\mathcal N}(C)\ge\tau\)
  \kw{then return} \(\operatorname{Serialize}(C)\)\\
\algline \kw{return} \(\mathsf{reject}\)\\
\end{algorithmic}
\end{algorithm}

The SFT input and label contain no element of \(\mathcal G\).  The parameter
values appear in Appendix~\ref{app:constants}.

\subsection{Complete RL Group-Scoring Algorithm}

The sampled policy is the same small open-weight model updated by RL.  Given
its rollouts, every reward component below is computed from program
executions, verifier outcomes, and deterministic clause provenance; no
external model enters the scoring function.

Executions of the masked loop yield reachable loop-head states \(\mathcal E\)
and target-independent negative traces \(\mathcal N\).  The lightweight
evaluator is three-valued:
\begin{equation}
  \operatorname{Eval}(c,s)\in
  \{\mathsf{true},\mathsf{false},\mathsf{unknown}\}.
  \label{eq:formal-state-eval}
\end{equation}
An unknown result leaves the negative trace unrejected and the candidate out
of the reachable-state front end's accepted set; the verifier handles syntax
and induction.
A trace counts as rejected only when at least one retained clause evaluates
definitively to false at one of its states:
\begin{equation}
  \operatorname{rej}_{\mathcal N}(C)
  =\{\nu\in\mathcal N:\exists s\in\nu,\exists c\in C,\
    \operatorname{Eval}(c,s)=\mathsf{false}\},
  \qquad
  \operatorname{cov}_{\mathcal N}(C)
  =\frac{|\operatorname{rej}_{\mathcal N}(C)|}{|\mathcal N|}.
  \label{eq:formal-coverage}
\end{equation}
Coverage is trajectory level and is defined only for
\(|\mathcal N|>0\).  Instances without retained negatives are rejected during
SFT construction.  For RL, the explicit fallback below uses a binary survivor
signal and never evaluates this ratio.

For an RL group \(y_{1:g_{\mathrm{rl}}}\), first form the same globally
deduplicated pool used at inference, and then invoke the verifier filter once:
\begin{align}
  A_i&=\operatorname{Admit}(y_i),
  &P^*&=\operatorname{Unique}\!\left(A_1\Vert\cdots\Vert A_{g_{\mathrm{rl}}}\right),
  \label{eq:formal-pooled-filter}\\
  C^*&=\operatorname{Filter}_{\mathcal L}(P^*),
  &V_i&=\left\langle c\in A_i\ \middle|\
    \operatorname{key}(c)\in\operatorname{keys}(C^*)\right\rangle,
  &R_i&=\operatorname{rej}_{\mathcal N}(V_i).
  \label{eq:formal-projection}
\end{align}
Here \(V_i\) assigns each pooled survivor back to every response that proposed
it.  A reachable-state front end evaluates the globally pooled candidates on
\(\mathcal E\) before the verifier processes the same pool.  If
\(f(\nu)=|\{j:\nu\in R_j\}|\), the reward is
\begin{equation}
  r_i=w_c\frac{|R_i|}{|\mathcal N|}
      +w_g\left(
        \frac{1}{|\mathcal N|}\sum_{\nu\in R_i}\frac{1}{f(\nu)}
      \right).
  \label{eq:formal-reward}
\end{equation}
Here \(w_c=1.0\) and \(w_g=0.3\) weight coverage and group credit.  The first term rewards each response by the negative coverage of
its pooled survivors.  The second is the response-level Shapley value of the
fixed post-filter coverage game
\(u(A)=|\bigcup_{i\in A}R_i|/|\mathcal N|\): repeated coverage is shared among
its proposers, while unique coverage receives full credit.  Thus
\(1/f(\nu)\) is a group-local difficulty weight: negatives rejected by many
responses receive less credit per response than rare negatives rejected by
few.  The negatives are nonuniformly weighted for each response, while their
credits still sum to one unit per covered trace across the group.
The per-response cap limits admission without penalizing
excess output.

\begin{algorithm}[H]
\caption{CRAFT reward for one RL group}
\label{alg:formal-rl-reward}
\begin{algorithmic}
\alginput{masked loop \(\mathcal L\), responses \(y_{1:g_{\mathrm{rl}}}\), reachable states \(\mathcal E\), negative traces \(\mathcal N\)}
\algoutput{response rewards \(r_{1:g_{\mathrm{rl}}}\)}
\algline \(A_i\leftarrow\operatorname{Admit}(y_i)\) for every \(i\)\\
\algline \(P^*\leftarrow\operatorname{Unique}(A_1\Vert\cdots\Vert
  A_{g_{\mathrm{rl}}})\)\cmnt{pool globally}\\
\algline \(C^*\leftarrow\operatorname{Filter}_{\mathcal L}(P^*)\)
  \cmnt{filter the pool once}\\
\algline \(V_i\leftarrow\langle c\in A_i\mid\operatorname{key}(c)\in
  \operatorname{keys}(C^*)\rangle\) for every \(i\)\cmnt{project by provenance}\\
\algline \kw{if} \(\mathcal N=\varnothing\) \kw{then}\\
\algline \algind \(r_i\leftarrow\mathbf 1[A_i\ne\langle\rangle\ \land\
  \operatorname{keys}(V_i)=\operatorname{keys}(A_i)]\) for every \(i\)\\
\algline \algind \kw{return} \(r_{1:g_{\mathrm{rl}}}\)\\
\algline \(R_i\leftarrow\operatorname{rej}_{\mathcal N}(V_i)\) for every \(i\)\\
\algline \(f(\nu)\leftarrow|\{j:\nu\in R_j\}|\) for every \(\nu\in\mathcal N\)\\
\algline \(r_i\leftarrow w_c|R_i|/|\mathcal N|+
  w_g|\mathcal N|^{-1}\sum_{\nu\in R_i}f(\nu)^{-1}\)
  for every \(i\)\\
\algline \kw{return} \(r_{1:g_{\mathrm{rl}}}\)\\
\end{algorithmic}
\end{algorithm}

A clause \(c\) \emph{semantically supports} \(d\) in \(D\) when
\(c\ne d\), \(\operatorname{Pres}_{\mathcal L,D}(d)\) holds, but
\(\operatorname{Pres}_{\mathcal L,D\setminus\{c\}}(d)\) does not.  The
implemented verifier exposes the broader, resource-dependent relation
\begin{equation}
\begin{split}
\operatorname{VSupport}_{\mathcal L,\delta}(c,d\mid D)
\iff{}&c\ne d\ \land
\operatorname{LoopVerify}_{\mathcal L,\delta}(D)[d]=\mathsf{valid}\ \land\\
&\operatorname{LoopVerify}_{\mathcal L,\delta}(D\setminus\{c\})[d]
  \ne\mathsf{valid}.
\end{split}
\label{eq:formal-verifier-support}
\end{equation}
This relation includes semantic support, generated-goal dependencies, and
resource effects.  Algorithm~\ref{alg:formal-rl-reward} can retain such a
clause through joint filtering.

\subsubsection{Shapley Credit: Definition and Equivalence}
\label{app:formal-shapley}

Fix the survivor bundles from one pooled filtering call and write
\(R_i=\operatorname{rej}_{\mathcal N}(V_i)\) and
\(f(\nu)=|\{j:\nu\in R_j\}|\).  For nonempty \(\mathcal N\),
\begin{equation}
  \begin{aligned}
  r_i^{\text{Clause-decomposed}}&=\frac{|R_i|}{|\mathcal N|},
  &\phi_i&=\frac{1}{|\mathcal N|}\sum_{\nu\in R_i}\frac{1}{f(\nu)},\\
  r_i^{\mathrm{Full}}&=r_i^{\text{Clause-decomposed}}+0.3\phi_i.
  \end{aligned}
  \label{eq:shapley-credit}
\end{equation}
\paragraph{Equivalence to the standard Shapley value.}
Equation~\eqref{eq:shapley-credit} gives a closed form for the standard
response-level Shapley value. For response \(i\), let \(V_i\)
be its projected bundle of clauses that survive filtering of the full response
group, and let \(R_i=\operatorname{rej}_{\mathcal N}(V_i)\).  A coalition
\(S\) contributes the projected bundles of its responses, whose rejected-trace
set is \(\bigcup_{i\in S}R_i\).

\begin{proposition}[Equivalence to standard Shapley]
\label{prop:shapley-closed-form}
Let \(G=\{1,\ldots,g\}\) index the responses above, and define their coverage
game by
\[
  u(S)=\frac{1}{|\mathcal N|}\left|\bigcup_{j\in S}R_j\right|,
  \qquad S\subseteq G.
\]
The standard Shapley value of response \(i\),
\begin{equation}
  \operatorname{Sh}_i(u)
  =\sum_{S\subseteq G\setminus\{i\}}
    \frac{|S|!\,(g-|S|-1)!}{g!}
    \bigl(u(S\cup\{i\})-u(S)\bigr),
  \label{eq:standard-shapley}
\end{equation}
is exactly
\[
  \operatorname{Sh}_i(u)
  =\frac{1}{|\mathcal N|}
    \sum_{\nu\in R_i}\frac{1}{f(\nu)}
  =\phi_i,
  \qquad
  f(\nu)=|\{j:\nu\in R_j\}|.
\]
Consequently, \(\sum_{i\in G}\phi_i=u(G)\).
\end{proposition}

\begin{proof}
The coefficient in Equation~\eqref{eq:standard-shapley} is the fraction of
permutations of \(G\) in which exactly the members of \(S\) precede \(i\):
there are \(|S|!\) orders before \(i\), \((g-|S|-1)!\) orders after it, and
\(g!\) orders in total.  Hence the standard subset definition equals the
expected marginal contribution of response \(i\) under a uniformly random
permutation.

For each trace \(\nu\), define
\(Q_\nu=\{j\in G:\nu\in R_j\}\).  If \(Q_\nu=\varnothing\), the trace never
contributes to the game.  Otherwise, it changes the coalition value exactly
once in any permutation: when the first response from \(Q_\nu\) is added.
By symmetry, each member of \(Q_\nu\) is first with probability
\(1/|Q_\nu|=1/f(\nu)\).  Thus trace \(\nu\) contributes
\(1/(|\mathcal N|f(\nu))\) to \(\operatorname{Sh}_i(u)\) when
\(\nu\in R_i\), and zero otherwise.  Summing these contributions over traces
gives Equation~\eqref{eq:shapley-credit}.  Finally, every trace covered by the
grand coalition allocates a total of
\(f(\nu)\cdot 1/(|\mathcal N|f(\nu))=1/|\mathcal N|\), proving efficiency.
\end{proof}

This equivalence is evaluated on the fixed response-level bundles
\(V_1,\ldots,V_g\) projected from the globally filtered survivor set.  It
therefore proves that the implemented allocation is the standard Shapley value
of this response-level coverage game. It assigns credit to responses using
fixed post-filter survivor bundles.


\subsection{Group-Relative Policy Update}
\label{app:grpo-details}

The rollout policy \(\pi_{\theta_{\mathrm{old}}}\) samples a group, which is
scored by Algorithm~\ref{alg:formal-rl-reward}.  Define
\begin{equation}
  \widehat A_i=
  \frac{r_i-\operatorname{mean}_{j}r_j}
       {\operatorname{std}_{j}r_j+\delta_{\mathrm{adv}}}.
  \label{eq:formal-advantage}
\end{equation}
For token \(a_{i,u}\) and its prefix
\(h_{i,u}=(\mathcal L,a_{i,<u})\), let
\begin{align}
  \rho_{i,u}(\theta)
  &=\frac{\pi_\theta(a_{i,u}\mid h_{i,u})}
          {\pi_{\theta_{\mathrm{old}}}(a_{i,u}\mid h_{i,u})},
  \label{eq:formal-token-ratio}\\
  q_{i,u}(\theta)
  &=\frac{\pi_{\mathrm{ref}}(a_{i,u}\mid h_{i,u})}
          {\pi_\theta(a_{i,u}\mid h_{i,u})},
  &k_{i,u}(\theta)&=q_{i,u}-\log q_{i,u}-1.
  \label{eq:formal-token-kl}
\end{align}
The nonnegative quantity \(k_{i,u}\) is the sampled conditional-token KL
surrogate used by the implementation.  Its expectation equals
\(D_{\mathrm{KL}}(\pi_\theta(\cdot\mid h_{i,u})\Vert
\pi_{\mathrm{ref}}(\cdot\mid h_{i,u}))\) when the token is sampled from
\(\pi_\theta(\cdot\mid h_{i,u})\).  During an update, however, the saved token
comes from \(\pi_{\theta_{\mathrm{old}}}\), so \(k_{i,u}\) is used as a
sampled penalty under the saved old-policy tokens.
The clipped group objective is
\begin{equation}
  \mathcal J(\theta)=\mathbb E\!\left[
  \frac{1}{g_{\mathrm{rl}}}\sum_i\frac{1}{T_i}\sum_{u=1}^{T_i}
  \left(
  \min\!\left\{\rho_{i,u}\widehat A_i,
  \operatorname{clip}(\rho_{i,u},1-\varepsilon,1+\varepsilon)\widehat A_i
  \right\}-\beta k_{i,u}\right)\right].
  \label{eq:formal-grpo}
\end{equation}
After optimizing this objective, the updated policy samples the next group.

\subsubsection{Group Reward Collapse and Shapley Tie-Breaking}
\label{app:formal-reward-collapse}

Here \emph{group reward collapse} means that all rewards in one sampled
GRPO group are equal.  Equation~\eqref{eq:formal-advantage} then gives zero
relative advantage to every response, so that group supplies no
reward-driven policy-gradient signal; a separate KL term can still act.
This condition concerns reward discrimination and does not require the
responses themselves to be identical.

If \(r_i^{\text{Clause-decomposed}}=c\) for every response in a group, then
\[
  r_i^{\mathrm{Full}}-\overline r^{\mathrm{Full}}
  =0.3(\phi_i-\overline\phi),\qquad
  \operatorname{Var}_i(r_i^{\mathrm{Full}})
  =0.09\operatorname{Var}_i(\phi_i).
\]
Thus Full restores a nonzero relative signal whenever equal-coverage
responses have different Shapley credits.  It can distinguish which
responses contribute less redundant coverage after Clause-decomposed has lost that
ranking signal.

For an illustrative calculation, let \(\mathcal N=\{a,b,c\}\) and three
responses reject \(R_1=\{a,b\}\), \(R_2=\{a,c\}\), and
\(R_3=\{a,c\}\), respectively.  All three Clause-decomposed rewards equal \(2/3\),
but the trace frequencies are \(f(a)=3\), \(f(b)=1\), and \(f(c)=2\).
Their Shapley credits are \((4/9,5/18,5/18)\), giving Full rewards
\((0.80,0.75,0.75)\).  Full therefore supplies positive relative credit to
the response that uniquely rejects \(b\).  This is a worked example of
the reward mechanism. It is illustrative and does not report a training outcome.

Full does not guarantee a noncollapsed group.  If every response rejects
the same trace set, both Clause-decomposed coverage and Shapley credit remain equal;
all-zero coverage also remains uninformative.  The mechanism can resolve
some coverage-score ties, but establishing a lower collapse frequency in
training requires logged groups from the two policies.

\subsection{Training--Inference Alignment}

Inference, SFT label construction, and RL group scoring all call the same
pooled map in Equation~\eqref{eq:formal-craft-map}; they differ only in how
the returned clauses are consumed.  Inference checks the restored goals, SFT
serializes an accepted fixed point, and RL projects the fixed point back to
the responses that proposed its surviving clauses.  Negative coverage supplies the target-independent
training surrogate; restored-goal verification supplies the final verdict.

\subsection{Extension to Verifiably Composable Generation}

\begin{definition}[Verifiably composable generation]
\label{def:vcg}
A task provides visible input \(x\), hidden goals \(\mathcal G\), an atomic
decomposition \(\operatorname{parts}_x\), a duplicate-removal operator
\(\operatorname{Unique}_x\), a subset filter \(\operatorname{Filter}_x\), a
composition operator \(\operatorname{Compose}_x\), and a validity predicate
\(\operatorname{Valid}_x\).  For responses \(y_{1:n}\), define
\begin{align}
  P_x(y_{1:n})
    &=\operatorname{Unique}_x\!\left(\bigcup_i
      \operatorname{parts}_x(y_i)\right),\\
  O_x(y_{1:n})
    &=\operatorname{Compose}_x\!\left(
      \operatorname{Filter}_x(P_x(y_{1:n}))\right).
  \label{eq:formal-general-pipeline}
\end{align}
The task is verifiably composable when the filter is conservative and sound:
it returns atoms from its input and
\(\operatorname{Valid}_x(O_x(y_{1:n}))\) holds whenever the filter succeeds.
It is target hidden when every operator above depends only on \(x\); the goals
\(\mathcal G\) are used only by a final task-specific check.  Complementary
composition occurs when \(O_x(y_{1:n})\) passes that final check although no
individual \(O_x(y_i)\) does.
\end{definition}

Loop-invariant generation instantiates the atoms with invariant clauses,
composition with conjunction, validity with joint inductiveness, and the
filter with bounded Houdini.  The transfer property for any other
instantiation satisfying Definition~\ref{def:vcg} is proved in
Appendix~\ref{app:composition-proofs}.

Definition~\ref{def:vcg} identifies the conditions needed to reuse CRAFT:
responses must decompose into reusable atoms, atoms must admit pooling, a
sound domain-specific filter must validate the pooled result, and provenance
must map accepted atoms back to responses.  Under these conditions, the same
inference and training construction applies after replacing the invariant
clauses, conjunction, and Houdini checks by the domain's atoms, composition
operator, and verifier.

For \emph{specification generation}, atoms can be precondition,
postcondition, frame, or invariant clauses; composition is conjunction; and
the filter checks syntax, consistency, and the corresponding proof
obligations.  Separate responses can contribute complementary clauses to one
verified contract.  For \emph{testcase generation}, atoms are executable
tests; composition is set union; the filter removes invalid or duplicate
tests; and structural or behavioral coverage supplies a target-independent
signal.  Separate responses can cover disjoint behaviors even when no one
response forms an adequate suite.

These domains share the formal interface and algorithms.  Each domain supplies
its verifier and measurable training signal.

\subsection{Theoretical Properties and Proofs}
\label{app:composition-proofs}

This subsection states and proves the properties of the algorithms defined
above.  Write
\(\mathcal L=(\mathsf{Init},B,S)\).

\subsubsection{Conjunction of Valid Invariants}

\begin{proposition}[Closure under conjunction]
\label{prop:conjunctive-closure}
If \(I_1,\ldots,I_m\) are valid invariants of the same loop, then
\(\bigwedge_j I_j\) is valid and is at least as strong as every \(I_j\).
\end{proposition}

\begin{proof}[Proof of Proposition~\ref{prop:conjunctive-closure}]
For each \(j\), validity gives \(\mathsf{Init}\Rightarrow I_j\); hence
\(\mathsf{Init}\Rightarrow\bigwedge_j I_j=I\).  Suppose \(I\wedge B\) holds before one
execution of \(S\).  Then \(I_j\wedge B\) holds for every \(j\), so preservation
of \(I_j\) establishes every \(I_j\) afterward.  Therefore \(I\) is preserved
and is a valid invariant.

Because a conjunction implies each conjunct, \(I\) is at least as strong as
every \(I_j\).
\end{proof}

\subsubsection{Mutual Support Among Arbitrarily Many Clauses}

\begin{proposition}[Mutual support]
\label{prop:mutual-support}
For a finite clause set \(C\), \(I(C)\) is valid if
\begin{equation}
  \mathsf{Init}\Rightarrow I(C),
  \qquad
  \forall c\in C:\ \operatorname{Pres}_{\mathcal L,C}(c).
  \label{eq:formal-mutual-support}
\end{equation}
No clause must be preserved in isolation.
\end{proposition}

\begin{proof}[Proof of Proposition~\ref{prop:mutual-support}]
The first premise is exactly establishment of \(I(C)\).  For preservation,
assume \(I(C)\wedge B\) before executing \(S\).  The second premise establishes
each \(c_j\) afterward.  Since every conjunct then holds in the same successor
state, their conjunction \(I(C)\) also holds.  Thus
\(\{I(C)\wedge B\}S\{I(C)\}\), and \(I(C)\) is valid.

No standalone preservation premise appears in this argument.  A clause may
use all other conjuncts in the pre-state, so clauses that fail preservation
separately may still be preserved jointly.  Establishment is different:
\(\mathsf{Init}\Rightarrow I(C)\) implies \(\mathsf{Init}\Rightarrow c_j\) for every \(j\).  Hence a
candidate false at some loop-entry state cannot be rescued by conjunction.
\end{proof}

\subsubsection{Greatest Jointly Inductive Candidate Subset}

Let \(C_0=C\), and let \(C_{j+1}\subseteq C_j\) be the set remaining after one
ideal Houdini deletion round in Equation~\eqref{eq:formal-houdini}.  Because \(C\) is
finite, this descending sequence reaches the fixed point
\(\mathsf H_{\mathcal L}(C)\).

\begin{theorem}[Candidate-relative optimality]
\label{thm:houdini-greatest}
With exact logical decisions, the Houdini fixed point
\(\mathsf H_{\mathcal L}(C)\) is jointly inductive and contains every jointly
inductive subset of \(C\).  Its conjunction is therefore the strongest
invariant obtainable by selecting clauses from this pool.
\end{theorem}
This candidate-relative property applies to ideal Houdini with the supplied
candidate pool. Bounded
filtering may discard a valid clause after timeout.  Joint checking permits
mutual support in preservation, but cannot rescue a clause false at loop entry.

\begin{proof}[Proof of Theorem~\ref{thm:houdini-greatest}]
At the fixed point, every \(c\in\mathsf H_{\mathcal L}(C)\) satisfies
establishment and preservation under the full conjunction
\(I(\mathsf H_{\mathcal L}(C))\).  Applying
Proposition~\ref{prop:mutual-support} shows that this conjunction is valid.

It remains to prove greatestness.  Let \(C'\subseteq C\) be any jointly
inductive subset.  We show by induction over deletion rounds that
\(C'\subseteq C_j\).  The base case \(C'\subseteq C_0=C\) is immediate.  Assume
\(C'\subseteq C_j\), and take any \(c\in C'\).  Joint establishment of \(C'\)
gives \(\mathsf{Init}\Rightarrow c\), so \(c\) cannot be deleted for establishment.  Joint
preservation gives

\[
  \{I(C')\wedge B\}\,S\,\{c\}.
\]

Since \(C'\subseteq C_j\), the antecedent \(I(C_j)\wedge B\) implies
\(I(C')\wedge B\).  Strengthening a valid Hoare precondition preserves
validity, so

\[
  \{I(C_j)\wedge B\}\,S\,\{c\}
\]

also holds.  Therefore no member of \(C'\) is deleted in round \(j\), proving
\(C'\subseteq C_{j+1}\).  At termination,
\(C'\subseteq\mathsf H_{\mathcal L}(C)\).  Because conjunction reverses set
inclusion, \(I(\mathsf H_{\mathcal L}(C))\) implies \(I(C')\).
\end{proof}

\paragraph{Pooling before filtering.}
Each separately filtered set \(\mathsf H_{\mathcal L}(A_i)\) is jointly
inductive.  By closure under conjunction, their union is a jointly inductive
subset of \(\bigcup_i A_i\).  Theorem~\ref{thm:houdini-greatest} therefore
gives Equation~\eqref{eq:main-pooling-inclusion}.  The inclusion can be strict
when a clause requires a supporting hypothesis from another response.
This ideal property does not ensure empirical monotonicity under finite
budgets: pooling changes proof obligations and may trigger timeouts
(Appendix~\ref{app:houdini-order}).

\subsubsection{Soundness of Resource-Bounded Houdini}

\begin{proposition}[Bounded-filter soundness]
\label{prop:bounded-houdini-sound}
Assume \(\operatorname{LoopVerify}_{\mathcal L,\delta}\) is sound.  If
\(\operatorname{Filter}_{\mathcal L}^{\delta,h_{\max}}(C)\) returns a proved
fixed point \(D\), then \(D\subseteq C\) and \(I(D)\) is inductive.  Timeout,
unknown, and exhaustion of the round budget cannot cause an unproved clause
set to be returned.
\end{proposition}

\begin{proof}
Every deletion round in Equation~\eqref{eq:formal-bounded-houdini-round}
selects a subset of its input, so \(D\subseteq C\).  At a returned fixed point,
every clause maps to \(\mathsf{valid}\) in that same family.  Verifier
soundness in Equation~\eqref{eq:formal-loop-verify-soundness} therefore gives
\(\operatorname{Ind}_{\mathcal L}(I(D))\).  A timeout or unknown removes the
affected clause, and reaching \(h_{\max}\) returns the empty result rather
than an intermediate set by Equation~\eqref{eq:formal-bounded-houdini}.
\end{proof}

\subsubsection{Supporting Clauses and Credit}

\begin{proposition}[Coverage credit omits pure support]
\label{prop:support-credit}
Suppose response \(i\) contributes a retained clause \(c\) that supports
another retained clause in the sense of Appendix~\ref{app:formal-method}, but
\(\operatorname{rej}_{\mathcal N}(V_i)=\varnothing\).  Then the coverage
term assigned to response \(i\) is zero, irrespective of whether removing
\(c\) would prevent the other clause from surviving.
\end{proposition}

\begin{proof}
The premise gives \(R_i=\varnothing\).  Therefore \(|R_i|=0\), so the
coverage term in Equation~\eqref{eq:formal-reward} is zero.
\end{proof}

This formal edge case has a negligible aggregate effect in our evaluation.
Post-filtering permits supporting clauses from different responses to survive
jointly. Pre-filtering removes clauses that cannot survive within their own
response. Under the longer verifier budget, their compose@\(k\)
difference ranges from \(-0.2\) to \(+0.6\) percentage points across the
reported budgets (Table~\ref{tab:prefilt-vs-postfilt}).  Thus, although the
post-filter coverage score does not credit pure indirect support, such
cross-response dependencies have negligible aggregate effect on verification
in this benchmark.

\subsubsection{Transfer to Verifiably Composable Problems}

\begin{proposition}[Composable transfer]
\label{prop:composable-transfer}
For any instantiation satisfying Definition~\ref{def:vcg}, the pooled output
\(O_x(y_{1:n})\) contains only atoms drawn from the responses and satisfies
\(\operatorname{Valid}_x\) whenever filtering succeeds.  Moreover, its
construction is independent of the hidden goals \(\mathcal G\).
\end{proposition}

\begin{proof}
Decomposition and union draw atoms only from \(y_{1:n}\);
\(\operatorname{Unique}_x\) removes candidates and the conservative filter
returns a subset, establishing provenance to the responses.  Filter soundness
gives \(\operatorname{Valid}_x(O_x(y_{1:n}))\).  Finally, every operator in
Equation~\eqref{eq:formal-general-pipeline} is a function of \(x\) and the
responses, not \(\mathcal G\), so hidden goals enter only the final check.
\end{proof}
